\section{Reye geometry, infinitesimal moduli, and the novelty boundary}
\label{sec:reye-priority}

We separate the classical moduli of webs and Reye congruences, the
exact differential calculation for the polarized-K3 map, and the two
nonreduced reconstruction theorems.  Generic web reconstruction uses
the deepest normal cone; it is not a consequence of the classical
nine-dimensional count.

\subsection{The classical nine-dimensional quotient}

For $e=4$, a relation space $R\subset\Sym^2V$ is a web of quadrics.
The bilinear incidence surface $Z_R$ of
\S\ref{sec:classical} is a K3 surface with its fixed-point-free
factor-exchange involution; its quotient is the classical Reye
Enriques surface associated with the web.  This construction and the
resulting congruence of lines are classical; see
\cite{Cossec,ArrondoSols,DolgachevReider,IK,DK,MMV}.

Arrondo--Sols compute that the Hilbert scheme of Reye congruences is
smooth of dimension $24$ and rational, and describe it by an open
Grassmannian parameter of the corresponding data.  After quotienting
by the $15$-dimensional projective coordinate group, the expected
moduli dimension is therefore nine.  In characteristic zero the
nodal Enriques locus is the classical nine-dimensional divisor in the
ten-dimensional Enriques moduli space; the modern theorem
\cite{MMV} identifies every classical nodal Enriques surface as a
Reye congruence.

Thus the equality
\[
 24-15=9
\]
is not a novelty claim of this paper.

\begin{proposition}[The classical tangent quotient]
\label{prop:reye-tangent}
At a general regular web with finite projective stabilizer, the tangent
space to the Reye moduli locus is the quotient of the $24$-dimensional
web tangent space by the $15$-dimensional infinitesimal
$\PGL(V)$-orbit:
\[
 T_{[R]}\mathcal M_{\mathrm{Reye}}
 \simeq
 T_R\Gr(4,\Sym^2V)/\mathfrak{pgl}(V),
 \qquad
 \dim T_{[R]}\mathcal M_{\mathrm{Reye}}=9.
\]
The differential of the polarized-K3 map of
Theorem~\ref{thm:main-moduli} factors through this quotient.
\end{proposition}
\begin{proof}
The first statement is the tangent-space form of the smooth
$24$-dimensional Hilbert calculation of Arrondo--Sols on the regular
web locus, followed by the infinitesimal quotient by projective
coordinates.  The stabilizer is finite at a general point, so the
orbit tangent has dimension $15$.

The construction $R\mapsto Z_R\simeq Y_R$ is equivariant for
$\PGL(V)$.  Passing from an embedded quartic to its polarized
isomorphism class kills the infinitesimal projective orbit; hence the
differential factors through the displayed quotient.
\end{proof}


\begin{proposition}[Web parameter and Reye Hilbert point]
\label{prop:web-hilbert-bridge}
Let
\[
 \mathcal G=\Gr(4,\Sym^2V)
\]
and let \(\mathcal G^{\mathrm{reg}}\) be the open on which the classical
Reye construction has the expected Hilbert polynomial.  The assignment
of a web to its Reye congruence is a \(\PGL(V)\)-equivariant morphism
\[
 \Phi:\mathcal G^{\mathrm{reg}}
 \longrightarrow \Hilb_{\mathrm{Reye}}.
\]
Under the Grassmannian parameterization of the regular Reye locus in
Arrondo--Sols, \(\Phi\) identifies a dense open of
\(\mathcal G^{\mathrm{reg}}\) with the corresponding
\(24\)-dimensional Hilbert component.  In particular, at a general
regular web,
\[
 d\Phi_R:
 \Hom(R,\Sym^2V/R)
 \xrightarrow{\ \sim\ }
 T_{[\operatorname{Reye}(R)]}\Hilb_{\mathrm{Reye}}.
\]
After quotienting by the infinitesimal projective orbit this gives
the tangent quotient in Proposition~\ref{prop:reye-tangent}.
\end{proposition}

\begin{proof}
Over \(\mathcal G\) the universal rank-four subbundle of
\(\Sym^2V\) gives the universal web.  The incidence equations defining
the Reye congruence are polynomial in the universal Pluecker
coordinates.  Restrict to the open on which their Hilbert polynomial
is the regular Reye polynomial.  Flatness there and the Hilbert
functor give the morphism \(\Phi\); equivariance is immediate from
the construction.

Arrondo--Sols describe the regular Reye Hilbert component by the open
Grassmannian parameter of the defining four-dimensional system of
quadrics.  With the present convention this parameter is precisely a
dense open of \(\Gr(4,\Sym^2V)\) (duality only changes whether one
writes the system or its annihilator).  Their smoothness and
dimension-\(24\) statement therefore identifies the tangent space of
the Hilbert point with the tangent space of this Grassmannian.  The
latter is \(\Hom(R,\Sym^2V/R)\), proving the displayed isomorphism.
Since the construction is \(\PGL(V)\)-equivariant, the orbit tangent
maps to the orbit tangent and the quotient gives
Proposition~\ref{prop:reye-tangent}.
\end{proof}

The role of Lemma~\ref{lem:differential} is now precise.  It is not
used to discover the number nine.  Its exact $25\times25$ integral
minor proves that the particular polarized K3 attached to our
multiplication problem does not collapse any of the $24$ embedded
web directions at the explicit rational point.  After projective
quotient, it proves rank nine for the map to polarized K3 moduli.
Since maximal rank is open, the induced map from the classical
nine-dimensional Reye deformation space to our polarized K3 image is
generically unramified.  This is the conceptual tangent statement
behind Theorem~\ref{thm:main-moduli}; the determinant is an exact
arithmetic witness for it.

\subsection{What the finite ambiguity means}

The preceding paragraph locates the Torelli ambiguity of the
polarized-K3 invariant alone.
The failure scheme first reconstructs the polarized K3
$(Y_R,\OO_{Y_R}(1))$.  A relation system additionally supplies the
free involution coming from factor exchange on the bilinear incidence
model and hence an Enriques quotient with its Reye data.  Classical
Reye geometry then has substantial rigidity.  In particular,
Dolgachev--Reider prove that for a Reye polarization on an Enriques
surface the stable rank-two bundle with
$c_1$ equal to that polarization and $c_2=3$ is the Reye bundle and
is unique up to isomorphism; it is the restriction of the universal
quotient bundle from $G(2,4)$.

Consequently the generic finite fibre in
Theorem~\ref{thm:main-moduli} should not be read as an uncontrolled
family of quadratic tensors.  Once a compatible Enriques quotient and
Reye polarization are fixed, the Grassmannian Reye bundle is rigid in
the preceding sense.  The remaining finite packet is concentrated in
the finite choices of compatible Reye/Enriques data over the recovered
polarized K3, together with the passage back to the web.  For the map defined by the polarized K3 alone, we prove finiteness
but do not assert a numerical degree or degree one; that enumeration
is distinct from generic reconstruction using the full failure
scheme.  The additional normal-cone invariant of
Theorem~\ref{thm:generic-intrinsic-web} separates this packet
without identifying its K3-only degree.

There is a second useful localization.  On the mixed-regular
projection-corank-two locus, the unsplit contact quartic of
\S\ref{sec:ranktwo-boundary} is the restriction of the already
recovered quartic $Y_R$ to the pencil of hyperplanes containing $H$.
Thus the generic corank-two contact cover itself is determined by the
polarized K3.  The following theorem turns the finite Torelli packet
into a precise finite-cover statement and identifies its monodromy.


\begin{theorem}[Generic Torelli packet as a finite etale monodromy cover]
\label{thm:torelli-monodromy}
Let \(\mathcal W\) be the irreducible nine-dimensional coarse
quotient of the regular web locus by projective equivalence on the
finite-stabilizer open, and let \(\mathcal K\) be the
nine-dimensional polarized-K3 image of
Theorem~\ref{thm:main-moduli}.  There are nonempty opens
\[
 \mathcal W^\circ\subset\mathcal W,\qquad
 \mathcal K^\circ\subset\mathcal K
\]
such that
\[
 \tau:\mathcal W^\circ\longrightarrow\mathcal K^\circ
\]
is finite etale of a constant degree
\[
 d_{\mathrm{Tor}}
 =
 [\,\C(\mathcal W):\C(\mathcal K)\,].
\]
Its geometric monodromy is transitive on the \(d_{\mathrm{Tor}}\)
points of a fibre.  Equivalently, if \(L\) is the normal closure of
\(\C(\mathcal W)/\C(\mathcal K)\), the monodromy group is the
permutation image of
\[
 \Gal(L/\C(\mathcal K))
\]
on the embeddings of \(\C(\mathcal W)\) in \(L\).

The degree-four contact incidence and its discriminant are pulled back
from the polarized quartic on \(\mathcal K^\circ\).  Hence they are
constant on a Torelli fibre and do not account for
\(d_{\mathrm{Tor}}\); any further separation of the finite packet
must use web-dependent structure beyond the polarized quartic.
Theorem~\ref{thm:generic-intrinsic-web} supplies such structure in
the deepest nonreduced normal cone, and
Corollary~\ref{cor:separate-torelli-packet} proves separation on a
nonempty open of this finite cover.
\end{theorem}

\begin{proof}
The source \(\mathcal W\) is irreducible because it is obtained from a
nonempty open of the irreducible Grassmannian
\(\Gr(4,\Sym^2V)\) by the projective-coordinate quotient on the
finite-stabilizer locus.  Theorem~\ref{thm:main-moduli} says that the
dominant map \(\tau:\mathcal W\dashrightarrow\mathcal K\) is
generically finite.  Thus the induced extension of function fields is
finite.  In characteristic zero it is separable.

Normalize a nonempty affine open of \(\mathcal K\) in
\(\C(\mathcal W)\).  After deleting the complement of the locus on
which the original rational map is defined, and then deleting the
branch and nonflat loci, one obtains a finite etale morphism from a
nonempty open of \(\mathcal W\).  Its degree is the function-field
degree displayed above.  Connectedness of the generic source, or
equivalently irreducibility of \(\mathcal W\), gives transitivity of
geometric monodromy.  The description by the normal closure is the
standard field-theoretic realization of that action.

The contact incidence is defined solely from the recovered embedded
quartic \(Y\subset\Pj^3\) by
\[
 (H,[\alpha])\quad\text{with}\quad H\subset\ker\alpha,
\]
and the discriminant is the discriminant of its restricted binary
quartic.  Both constructions commute with isomorphism of the polarized
quartic.  They therefore descend to \(\mathcal K^\circ\) and pull
back along \(\tau\), proving the last assertion.
\end{proof}



\subsection{Comparison with nonreduced reconstruction theory}

The use of nilpotent layers has classical antecedents.  A ribbon is a
square-zero double structure with an invertible conormal module, and
Bayer--Eisenbud developed its intrinsic and canonical geometry
\cite{BayerEisenbud}.  Manolache surveys Cohen--Macaulay nilpotent
structures and their filtrations \cite{Manolache}.  Primitive multiple
schemes carry a canonical filtration by lower multiplicities in the
sense of Drézet \cite{Drezet}.

Our local corank-one model
\[
 t(t,f)=(t)\cap(t,f)^2
\]
indeed has a square-zero nilradical line on the embedded support, and
the corank-two model has one further nonzero nilpotent power.  We do
not claim that the existence of a nilpotent filtration is new.  The
specific contribution is the following chain of implications, which
does not occur in the cited multiple-structure theory:
\[
 \text{multiplication Fitting scheme}
 \Longrightarrow
 \bigl(\ZZ_1,\widehat\NN|_{\ZZ_1}\bigr)
 \Longrightarrow
 \omega_{\ZZ_1}\otimes\widehat\NN^{-9}
 \Longrightarrow
 (Y_R,\OO_{Y_R}(1)),
\]
together with
\[
 \ZZ_2=\text{rank-two Schubert support}
\]
on the mixed-regular higher-corank locus.  The ambient failure scheme
is not assumed to be a primitive multiple scheme, and the filtration
is defined by annihilators of powers of its intrinsic nilradical
rather than by choosing an ambient thickening.

\subsection{Ballico's failure-locus framework}

The complete theorem text of Ballico's 1993 article
\cite{Ballico93} has not been obtained for theorem-level comparison.  We therefore do
not infer either anticipation or non-anticipation from its title or
metadata.  A theorem-level comparison is possible with Ballico's
accessible 1996 continuation \cite{Ballico96}.  There the basic data
are an integral projective variety $X$, a very ample line bundle $L$,
the complete embedding by $H^0(X,L)$, and multiplication maps such as
\[
 H^0(L)\otimes H^0(L^k)\longrightarrow H^0(L^{k+1}).
\]
The principal results propagate failure of quadratic or higher
normality to suitable linear sections and finite subschemes; the
parameter Grassmannian is a Grassmannian of projective linear
sections.

The present construction has a different input and a different
scheme-theoretic output.  The fixed object is the finite local
cube-zero algebra $B_R$; a point of $X_R$ is a unital subspace
$\C1\oplus K\subset B_R$; and the failure scheme is the zeroth Fitting
scheme of
\[
 \Sym^m(\OO\oplus\mathcal K)\longrightarrow B_R\otimes\OO.
\]
Its embedded associated supports, nilradical powers, and the
polarized K3 reconstructed from their line bundles are part of the
statement.  Ballico's 1996 theorems are therefore an important
precedent for organizing multiplication failure over a Grassmannian,
but they do not by themselves supply the Fitting normal forms or the
nilpotent-depth reconstruction proved here.  This comparison is
deliberately limited to the theorem text actually inspected; the
separate 1993 comparison remains documentary work rather than a
premise in a novelty claim.

\subsection{Residue of the main theorem}

After these classical comparisons, the claims that remain specific to
this paper are concise.

First, an abstract multiplication-failure scheme reconstructs the
quartic polarization by the intrinsic line
$\omega_{E_R}\otimes\NN_R^{-9}$ and its section ring.  Second, the
same nonreduced scheme has a functorial annihilator filtration whose
next nonzero stage recovers the first higher-corank Schubert support.
Third, on the mixed-regular rank-two locus the unsplit contact ideal
has a complete collision specialization law, the fifth-power vertex
component remains genuinely primary, and the associated graded
nilpotent algebra identifies both the contact cone and the rank-two
Schubert support.  The first decomposable mixed-kernel wall is also
computed scheme-theoretically.  Finally,
Theorem~\ref{thm:torelli-monodromy} identifies the inverse problem
for the K3 alone with a finite etale monodromy cover, while
Theorem~\ref{thm:generic-intrinsic-web} proves that the actual
nonreduced failure cone generically recovers the web and separates
that packet.  These statements are
independent of whether the nine-dimensional Reye family itself is
classical.
